%!TEX root = ../template.tex
%%%%%%%%%%%%%%%%%%%%%%%%%%%%%%%%%%%%%%%%%%%%%%%%%%%%%%%%%%%%%%%%%%%%
%% chapter4.tex
%% NOVA thesis document file
%%
%% Chapter with Implementation
%%%%%%%%%%%%%%%%%%%%%%%%%%%%%%%%%%%%%%%%%%%%%%%%%%%%%%%%%%%%%%%%%%%%


\typeout{NT FILE chapter4.tex}%

\chapter{Implementation}
\label{cha:implementation}

This section is split into different sections explaining the project's developed software.

%\section{Prototype Hardware}
%\label{sec:prototypeHardware}

%Two prototypes were developed and they are composed by an Arduino 33 nano \gls{BLE}  Rev 2, an \gls{EMG} and a small power bank to provide power (see \ref{sec:prototype}).

%Foto do sistema completo

%\subsection{Arduino}
%\label{ssec:arduino}

%As mentioned in the introduction of this chapter, the Arduinos used are two Arduino 33 nano \gls{BLE} Rev 2 which is a 3.3V board of only 45x18mm of area. It features a powerful processor, the nRF52840 from Nordic Semiconductors, a 32-bit 
%ARM® Cortex®-M4 CPU running at 64 MHz. In terms of sensors, the Arduino has a 9 axis \gls{IMU} meaning it includes an accelerometer, a gyroscope, and a magnetometer with 3-axis resolution each. As its name mentions it also has \gls{BLE} 
%communication which is vital for this project~\cite{arduinoDoc}.



\section{Arduino Code}
\label{sec:arduinoCode}

The Arduino programming language is a modified C/C++ language that simplifies the board coding that is run on the Arduino's IDE.
Initially there were two code files for the Arduinos, one for the sensors calibration (specifically the magnetometer) and the other one is run for reading and sending the data. The calibration 
was eventually integrated in the main program and only a single program is needed.


\subsection{Calibration}
\label{ssec:calibration}

To ensure accurate and reliable measurements, the magnetometer and the gyroscope undergo a calibration process. Since the magnetometer's values depend heavily on the surrounding environment, a calibration routine must be executed 
before using the device. During this process, the program records the magnetometer's minimum and maximum values while the device is moved following a continuous motion resembling an infinity 
symbol ($\infty$) and rotated around its axis until the user decides it's enough.

The offset for each axis is then determined by the following equation:

\begin{equation}
  offset = minValue + (maxValue - minValue) / 2
  \label{eq:offsetEQ}
\end{equation}

And the scale is calculated through the next equation:

\begin{equation}
  scale = (maxValue - minValue) / 2
  \label{eq:scaleEQ}
\end{equation}

In these expressions the minValue and the maxValue correspond to the lowest and highest recorded values for each axis. The corrected measurements for each axis can then be expressed as:

\begin{equation}
  value = (rawValue - offset) / scale
  \label{eq:magnetometerValueEQ}
\end{equation}

The gyroscope also requires calibration to eliminate zero-rate bias (drift) that occurs even when the sensor is stationary. This process consists on leaving the device completely still for 5 seconds while the raw angular velocity 
is measured and stored. At the end, the average of these readings for each axis is the offset used in the main program:

\begin{equation}
  value = rawValue - offset
  \label{eq:gyroscopeValueEQ}
\end{equation}

After connecting both devices, in the Data Menu (\autoref{sssec:data_menu}), there is a button that takes us to the Calibration Screen (\autoref{sssec:calib_menu}) where the user can proceed with that process. This should 
always be done before collecting data in a new location or time. In case of a disconnection, when the device reconnects, the app automatically loads the last saves calibration values.
When the arduino receives the calibration command (either gyroscope or magnetometer), it starts the respective operation until, in the magnetometer's case, the user ends it or, in the gyroscope case, it finishes automatically. Afterwards 
the values are set as offsets and they're sent to the app where they are saved for later upload.


\subsection{Sensors' Data Readings And Feature Calculations}
\label{ssec:arduino_sensors_data_read}

To achieve consistent gait analysis and movement recognition, the system requires a high frequency of \gls{EMG} and \gls{IMU} data collection. However \gls{BLE} has strict restrictions on the packet output frequency specially since we're 
dealing with two devices at high frequencies (they may collide with each other).

\subsubsection{Sampling and Windowing}
\label{sssec:sampling_windowing}

Raw \gls{EMG} signals contain significant noise and offset so, before feature extraction, a High-Pass Filter is applied with a cut-off frequency of approximately 20 Hz ($\alpha$ =0.888). The \gls{IMU} data is fed to the \gls{ML} trainings raw.
The microcontrollers operate in a non-blocking loop architecture designed to handle both the \gls{EMG} and \gls{IMU} acquisitions and different rates:

\begin{itemize}
  \item High-Frequency \gls{EMG} Readings (1000 Hz) - To obtain relevant features from the sensor it is read every 1 ms.
  \item Kinematic (\gls{IMU}) \& Transmission (50 Hz) - Every 20 ms the \gls{IMU} data is measured right before the ble communication since no calculations are done to that data.
\end{itemize}

The coordination of both timings is managed through a windowing approach. A sliding window of 20 \gls{EMG} samples is filled and once full (after 20 ms), the system triggers an event:

\begin{enumerate}
  \item The \gls{IMU} sensors are read.
  \item The time-domain features are calculated from the last 200 \gls{EMG} samples (see \autoref{sssec:emg_features_calculation})
  \item All the data is bundled in a packet and waits to be sent (see \autoref{sssec:arduino_payload}) %mencionar próxima secção
\end{enumerate}


\subsubsection{\gls{EMG} Features Calculation}
\label{sssec:emg_features_calculation}

As mentioned before, \gls{EMG} signals are non-stationary and contain significant noise making them unsuitable for direct raw transmission via \gls{BLE} due to bandwidth constraints. Furthermore, many studies indicate that raw \gls{EMG} 
underperforms on a \gls{ML} classifier environment when compared against \gls{EMG} features.
The microcontroller extracts time-domain features since they are computationally more efficient than the frequency-domain-features that have the heavy processing load of the Fast Fourier Transforms (FFT).
Each calculation processes a window of 200 samples (200 ms) resulting in six measurements:

\begin{enumerate}
  \item \gls{RMS} \autoref{sssec:root_mean_square}
  \item \gls{MAV} \autoref{sssec:mean_absolute_value}
  \item \gls{VAR} \autoref{sssec:variance}
  \item \gls{ZC} \autoref{sssec:zero_crossing}
  \item \gls{SSC} \autoref{sssec:slope_sign_changes}
  \item \gls{WL} \autoref{sssec:waveform_length}
\end{enumerate}

%Talvez um examplo de um sinal emg e o significado de cada feature

\subsection{Bluetooth Low Energy}
\label{ssec:arduino_ble}

For the Bluetooth connection we first have to assign a name to the device and a unique address for its service and characteristics. In this project three characteristic are used.
Then, as explained in section X, we advertise the service to the nearby devices and when there is a connection the last saved calibration values are uploaded and then the data collection is started and 
the characteristic is updated.%falta número da secção

\subsubsection{GATT Service}
\label{sssec:arduino_gatt}

The system defines a unique service identified by a 128-bit UUID. Within this service, three specific Characteristics are utilized to handle data streaming and bi-directional control command flow.

\begin{table}[ht]
  \centering
  \caption{Bluetooth Characteristic Configuration}
  \label{tab:gatt_config}
  \renewcommand{\arraystretch}{1.5}
  \begin{tabular}{|m{3.2cm}||m{2.8cm}|m{2.2cm}|m{5cm}|}
    \hline
    \multicolumn{4}{|c|}{\textbf{Characteristic Configuration}} \\
    \hline
    Characteristic Names&Properties&UUID Prefix&Function\\
    \hline
    Data Stream& READ \& \newline NOTIFY&1b4867e2&The android app listens to this characteristic and is notified whenever the devices upload data\\
    Calibration \newline Command& WRITE&1b4867e3&When the user starts calibrating the android app sends commands via this characteristic\\
    Calibration Data& READ \& \newline NOTIFY&1b4867e4&Here the devices upload the progress of the calibrations and the final results\\
    \hline
  \end{tabular}
\end{table}

\subsubsection{Payload Optimization}
\label{sssec:arduino_payload}

 A standard 32-bit floating-point representation would require 60 bytes per sample, which is inefficient for BLE bandwidth so in order to reach the objective of 50 Hz transmission rate for the nine IMU axis and six EMG features the 
 payload has to be highly optimized. With this in mind, all the float sensor values were converted to 16 bit signed integers reducing the packets significantly while maintaining enough precision for analysis.

 Each sensor frame consists of:

\begin{enumerate}
  \item Accelerometer Data (6 bytes): X, Y, Z ($\times100$ scaling)
  \item Gyroscope Data (6 bytes): X, Y, Z ($\times100$ scaling)
  \item Magnetometer Data (6 bytes): X, Y, Z ($\times100$ scaling)
  \item \gls{EMG} Features (12 bytes): gls{RMS}, gls{WL}, gls{ZC}, gls{MAV}, gls{SSC}, gls{VAR} (Mixed scaling: $\times10000, \times10000 or \times10$)
\end{enumerate}

\begin{equation}
  Total Size = 18 (\gls{IMU}) + 12 (\gls{EMG}) = 30 Bytes
  \label{eq:packetSize}
\end{equation}

Although the packet size has been reduced, transmitting an individual packet every 20 ms from two separate device still introduces a \gls{BLE} protocol overload. To get around this problem, the packets are bundled in groups of 5 ($30 bytes\times5 packets=150 bytes$).
This strategy effectively utilizes the large MTU (Maximum Transmission Unit) negotiated with the Android device, reducing the communication frequency to exactly once every 100 ms (for each device) reducing cpu usage and radio congestion. 
To make the system even more robust, a FIFO (First In First Out) buffer with a capacity of 150 packets was implemented. Whenever the communication with the Android device fails, that data is then saved in the buffer and the device 
tries to send it again later.
One other challenge introduced by the \gls{BLE} transmission and buffering implemented is the blocking nature of the communication operations. Pushing a 150 Byte bundle to the Bluetooth stack will momentarily block the main processing loop. 
While the microprocessor waits for the radio to finish, it will miss some of the 1 ms window defined for the \gls{EMG} sampling leading to gaps in the timeline. To circumvent this, the system tracks the elapsed time strictly via a 
microsecond hardware timer. Immediately before constructing and transmitting the next \gls{BLE} bundle, the firmware compares the current time against the scheduled \gls{EMG} sampling interval. If the previous transmission caused a delay, 
the system enters a rapid catch-up loop, executing the analog read and filtering functions multiple times sequentially until the "time debt" is cleared ensuring that while the transmission may introduce some jitter in the packet arrival time 
at the Android device, the sampling period remains consistent within the microcontroller's memory, preserving the temporal integrity required for accurate frequency-domain features.


\section{Android Application}
\label{sec:android_app}

%diagrama de blocos

As mentioned in the Work Proposal (\autoref{cha:workProposal}), a dedicated Android application was developed in Kotlin, to serve as a bridge between the prototypes and the user/database. This programming 
language is used to interoperate fully with Java developed by JetBrains and since 2019 it has been the preferred language by Google for Android app development.
The application handles the \gls{BLE} connection with the devices and provides a user interface for researchers, athletes and coaches to visualize the collected data and the model's predictions in real time.

The following section provides a detailed description of the software architecture of the application. Each file is explained according to its function within the system, with emphasis on how data flows from the BLE layer to the 
machine learning classifier and finally to the user interface. This breakdown clarifies the responsibilities of each component, highlights the interaction between modules, and demonstrates how the application was structured to 
ensure modularity, maintainability, and extensibility.

This application follows a modular architecture separating responsibilities between the user interface, data handling, and Bluetooth communication. At its core, the program follow the MVVM (Model-View-ViewModel) design pattern, improving 
testability, maintainability and division of tasks.

The MainActivity.kt is the entry point of the application, initializing the navigation host and providing the top-level container for the user interface. From this activity, the navigation system manages the transitions between the 
windows (implemented as Jetpack Compose composables), each of which corresponds to a specific functionality of the program.
The Android application serves as bridge between the devices, the user and the database. 



\subsection{User Interface}
\label{ssec:android_interface}

The application consists in 4 different screens that define the layout and interactive elements that the user sees. Each one of these windows has a corresponding ViewModel that acts as a connection point between the user interface 
and the application logic. These ViewModels handle any calls made by the screens and any changes to relevant operations such as the bluetooth connection.

At first, when the app is launched, some necessary permissions are checked, specifically the Bluetooth (\texttt{BLUETOOTH\_SCAN and BLUETOOTH\_CONNECT}), 
the location (\texttt{ACCESS\_FINE\_LOCATION and ACCESS\_COARSE\_LOCATION}) and the internet access permissions. Even though the app doesn't use the user's location, its permission is required for \gls{BLE} in Android 12 and higher. 

\begin{figure}[htbp]
  \centering
  \includegraphics[width=0.8\linewidth]{bleTurnOn.jpg}
  \caption{Request message to turn on Bluetooth}
  \label{fig:turnOnBluetooth}
\end{figure}

\subsubsection{Start Menu}
\label{sssec:start_menu}

%falar de cada viewmodel
This is the first menu shown to the user and it has the simple task of guiding the user through the application's functions. It has one button that takes us to the profile menu (\ref{sssec:profile_menu}) and another that takes us to 
data menu(\autoref{sssec:data_menu}).

\begin{figure}[htbp]
  \centering
  \includegraphics[width=0.8\linewidth]{startScreen.jpg}
  \caption{Start Menu}
  \label{fig:startScreen}
\end{figure}

\subsubsection{Profile Menu}
\label{sssec:profile_menu}

In this menu every saved user profile is listed and the user is able to create new profiles and delete existing ones. When creating a profile the app generates a unique id and the user is asked for a name and height. 
The list of users is stored locally and is available in every session.

\begin{figure}[ht]
    \centering
    \begin{subfigure}[b]{0.5\textwidth}
        \centering
        \includegraphics[width=\textwidth]{usersScreen.jpg}
        \caption{Profiles' Menu}
        \label{fig:usersScreen}
    \end{subfigure}
    \hfill 
    \begin{subfigure}[b]{0.35\textwidth}
        \centering
        \includegraphics[width=\textwidth]{newUserWindowOnly.jpg}
        \caption{New User Window}
        \label{fig:newUserWindow}
    \end{subfigure}
    
    \caption{Profile's Menu and User Creation}
    \label{fig:dataScreens}
\end{figure}


\subsubsection{Data Menu}
\label{sssec:data_menu}

When this menu is opened, the app tries to connect to both devices and when able, starts exchanging data. The data is then shown to the user. On the left the user can select an active user profile and a label for the current movement. On the right, 
from top to bottom, the \gls{ML} models' predictions are shown (\autoref{}). Beneath the models' projections, there are 2 switches - one to start/stop the \gls{ML} server connection and another one to start/stop sending the collected data to the 
InfluxDB database. At last there is a button that takes the application to the Calibration Menu (\autoref{sssec:calib_menu}).

%falar do VariabResult

\begin{figure}[ht]
    \centering
    \begin{subfigure}[b]{0.45\textwidth}
        \centering
        \includegraphics[width=\textwidth]{dataScreenNoCon.jpg}
        \caption{Without devices' connections}
        \label{fig:dataScreenNoC}
    \end{subfigure}
    \hfill 
    \begin{subfigure}[b]{0.45\textwidth}
        \centering
        \includegraphics[width=\textwidth]{dataScreenWithCon.jpg}
        \caption{When both devices are connected}
        \label{fig:dataScreenWithCon}
    \end{subfigure}
    
    \caption{Data Menu}
    \label{fig:dataScreen}
\end{figure}


\subsubsection{Calibration Menu}
\label{sssec:calib_menu}

This window provides a dedicated interface for managing the sensor calibration routines which are critical for ensuring the accuracy of the \gls{IMU}. As mentioned in the Arduino section (\autoref{ssec:calibration}), due to the inherent hardware 
drift and environmental magnetic interference, the gyroscope and magnetometer have to be recalibrated periodically for better and more accurate results.
In this page the user can trigger the calibration methods that send specific commands to the devices with the instructions needed. For the gyroscope, the firmware executes a static calibration by averaging 1000 samples to compute and nullify 
the offset. For the magnetometer, the system performs a dynamic hard-iron calibration, recording the maximum and minimum magnetic field intensities during a multi-axis rotation to calculate precise spatial offsets and scaling factors.

\begin{figure}[htbp]
  \centering
  \includegraphics[width=0.8\linewidth]{calibScreen.jpg}
  \caption{Calibration Menu}
  \label{fig:calibMenu}
\end{figure}

The gyroscope calibration for both devices is done at the same time while the magnetometers' are executed one at a time. For each operation the app has a respective button and while ongoing calibration, the app display some real-time progress 
to the user via a dedicated \gls{BLE} characteristic.

\begin{figure}[ht]
    \centering
    \begin{subfigure}[b]{0.45\textwidth}
        \centering
        \includegraphics[width=\textwidth]{calibScreenIMU.jpg}
        \caption{Gyroscope Calibration in Progress}
        \label{fig:calibMenuIMU}
    \end{subfigure}
    \hfill 
    \begin{subfigure}[b]{0.45\textwidth}
        \centering
        \includegraphics[width=\textwidth]{calibScreenFinished.jpg}
        \caption{Calibration Finished}
        \label{fig:calibMenuFinished}
    \end{subfigure}
    
    \caption{Calibration Menu}
    \label{fig:calibMenu2}
\end{figure}

When all the calibration operations are finished, the following screen (\autoref{fig:calibMenuFinished}) is presented and the values are saved to be uploaded when the devices reconnect


\subsection{Arduino BLE Communication}
\label{ssec:android_com_arduino}

In Android \gls{BLE}  data is handle asynchronously which means that the app receives callbacks which have to be handled. For the \gls{BLE} communication Android provides a package with a layered architecture:

\begin{itemize}
  \item BluetoothAdapter - Where every bluetooth connection starts. It allows scanning for nearby available devices and initiating connections.
  \item BluetoothDevice - Represents another \gls{BLE}  device with one or more services.
  \item Bluetooth\gls{GATT} - It's the protocol for \gls{BLE}  services and characteristics. It manages the connection lifecycle and provides callbacks for asynchronous communication.
  \item Bluetooth\gls{GATT} Service - Each service has a group of associated characteristics and a specific function (e.g., location service)
  \item Bluetooth\gls{GATT} Characteristic - It's the smallest unit of data in \gls{BLE} BLE and it can be read, written or subscribed to (to ge notified whenever it changes) depending on its attributes.
\end{itemize}

Since the are two devices, each with its own uuid (see \autoref{ssec:arduino_ble}), each device has its own manager that handle its callbacks (VariabBLEReceiveManager.kt) and these two managers are handled by the DualBLEReceiveManager 
that coordinates them. 


\subsubsection{Receive Managers}
\label{sssec:manager}

Each receive manager is responsible for managing the lifecycle of the corresponding \gls{BLE} connection and handling every callback. Depending on the uuid of the connected device, the manager knows with which device it is working with 
(if left or right). 

First, using the BluetoothAdapter (see \autoref{ssec:android_com_arduino}), the manager searches for the respective device using its name and uuid and once discovered a \gls{GATT} connection is tried. After establishing the connection, 
the manager finds the device's services. If a device disconnects the application cleans up its resources and starts trying to reconnect.

In the first approach the system was designed with separate \gls{GATT} characteristics for different sensor axes and \gls{EMG} data. However, relying on multiple simultaneous asynchronous notifications caused the Android \gls{BLE} stack to 
frequently drop packets and ignore callbacks. To resolve this without resorting to an inefficient sequential polling queue, the system was optimized by bundling all sensor data into a single 150 byte payload transmitted via a single 
characteristic.

The managers subscribe to the data characteristic of the respective device and rely on the onCharacteristicChanged callback. This allows the microcontrollers to asynchronously push data bundles at 10 Hz. Upon receiving the payload, the 
managers decode the byte arrays into their respective float and integer values. If the data upload switch is active, this payload is appended with the active userId, the selected movement label, and a timestamp, and is immediately sent 
to the database.


\subsection{Communication with InfluxDB}
\label{ssec:android_com_data_base}

As mentioned before the application uses InfluxDB as its database. The interface between the application and the database is implemented in two files, one containing the function that sends data and the other one has the sensitive 
client/database info (token, database url and organization and bucket name). This is a coroutine-based execution meaning that it will never stop or delay the main user interface thread and that it can be run multiple times in parallel 
without affecting the application performance.

Since InfluxDB organizes data using measurements, tags, fields, and timestamps, the application must first parse the incoming \gls{BLE} payloads into the InfluxDB Line Protocol. This is a text-based format that optimizes ingestion 
speed by representing the data points in a single string structure:

\begin{itemize}
    \item Measurement - movement\_data
    \item Tags - sensor(left or right) and userId
    \item Fields (data values) - accelerometer, gyroscope, magnetometer and \gls{EMG} data
    \item Timestamp - Instant.now() (registers that precise moment)
\end{itemize}

Once formatted, the data is transmitted to the cloud via an asynchronous HTTP POST request to the InfluxDB API, ensuring any network timeouts or errors are caught and handled without freezing the application's user interface.

\subsection{User Database}
\label{ssec:android_users}

To track the origin of the physiological signals, the Android application has a robust local profiling system. Before data collection begins, the user should identify the active subject (otherwise the app send a null for the id to the database). 
This is managed locally on the Android device using the Android Room Persistence Library, which provides an abstraction layer over an SQLite database to allow fluent database access while harnessing the full power of native \gls{SQL}.

\subsubsection{User Data Model and Data Access Object}
\label{sssec:android_users_model}

The database system is defined by a data entity (user.kt) that represents a subject. Each subject is assigned an auto-generated primary key (id), a name, and height in centimetres. 
Interactions with the database are strictly mediated through a Data Access Object (DAO). The UserDao defines the \gls{SQL} operations required to insert, delete, and retrieve subject profiles. To maintain a fluid user interface and prevent 
blocking the application's main thread during disk I/O operations, all database queries are executed asynchronously using Kotlin Coroutines (suspend functions).

\subsubsection{Repository Pattern and Reactive State}
\label{sssec:android_users_repository}

To separate the data persistence logic from the User Interface, the architecture employs the Repository Pattern (UserRepository.kt). The repository acts as the single source for user data across the entire application.
It abstracts the DAO layer, providing clean functions for the application's ViewModels to fetch lists of available users or create new ones while maintaining recollection of the active user.

The active user's ID is exposed to the rest of the application as a reactive StateFlow and is sent along side the sensors' data to the database. This modern reactive programming paradigm ensures that whenever the active user changes in the 
User Menu, any observing component in the application (such as the \gls{BLE} transmission service) is instantly and automatically updated with the new context.

\subsection{Communication with the Machine Learning server}
\label{ssec:com_ml_server}






\section{Machine Learning Code}
\label{sec:ml_code}



\subsection{Models' Training}
\label{ssec:ml_training}


\subsection{Random Forest}
\label{ssec:ml_rf}

\subsection{Tiny \gls{ML}}
\label{ssec:ml_tml}